\subsection{滤子基的像与逆像}

设 \(\mathfrak{B}\) 是集合 \(\mathbf{X}\) 上的滤子基，并设 \(f\) 是 \(\mathbf{X}\) 到集合 \(\mathbf{X}'\) 中的映射；则 \(f(\mathfrak{B})\) 是 \(\mathbf{X}'\) 上的滤子基，因为由关系 \(\mathbf{M} \neq \emptyset\) 可得 \(f(\mathbf{M}) \neq \emptyset\)，且我们有 \(f(\mathbf{M} \cap \mathbf{N}) \subset f(\mathbf{M}) \cap f(\mathbf{N})\)。若 \(\mathfrak{B}_1\) 是细于以 \(\mathfrak{B}\) 为基之滤子的某个滤子的基，则 \(f(\mathfrak{B}_1)\) 是细于以 \(f(\mathfrak{B})\) 为基之滤子的某个滤子的基（第 3 目，命题 4）。

命题 10. 若 \(\mathfrak{B}\) 是集合 X 上的超滤子基，且 \(f\) 是 X 到集合 \(\mathbf{X}'\) 中的映射，则 \(f(\mathfrak{B})\) 是 \(\mathbf{X}'\) 上的超滤子基。

设 \(\mathbf{M}'\) 是 \(\mathbf{X}'\) 的子集。若 \(\overline{f}^{1}(\mathbf{M}')\) 包含某个集合 \(\mathbf{M}\)（属于 \(\mathfrak{B}\)），则 \(\mathbf{M}'\) 包含 \(f(\mathbf{M})\)；若不然，则 \(\complement \overline{f}^{1}(\mathbf{M}') = \overline{f}^{1}(\complement \mathbf{M}')\) 包含某个集合 \(\mathbf{N}\)（\(\mathfrak{B}\) 中的）（第 4 目，命题 5），从而 \(\complement \mathbf{M}'\) 包含 \(f(\mathbf{N})\)。于是结论由第 4 目命题 6 得出。

特别考虑 \(f\) 是典范单射 \(\mathbf{A} \to \mathbf{X}\)（\(\mathbf{A}\) 为集合 \(\mathbf{X}\) 的子集）的情形。若 \(\mathfrak{B}\) 是 \(\mathbf{A}\) 上的滤子基，则 \(f(\mathfrak{B})\) 是 \(\mathbf{X}\) 上的滤子基。滤子 \(\mathfrak{F}\)（\(\mathbf{X}\) 上的，由 \(f(\mathfrak{B})\) 生成的）称为由 \(\mathfrak{B}\)（把 \(\mathfrak{B}\) 视为 \(\mathbf{X}\) 上的滤子基时）生成的滤子。若 \(\mathfrak{B}\) 是 \(\mathbf{A}\) 上的超滤子基，则由命题 10，它也是 \(\mathbf{X}\) 上的超滤子基。

接下来考察滤子基的逆像是否为滤子基。设 \(\mathfrak{B}'\) 是集合 \(\mathbf{X}'\) 上的滤子基，并设 \(f\) 是集合 \(\mathbf{X}\) 到 \(\mathbf{X}'\) 中的映射；则 \(\overline{f}^{1}(\mathfrak{B}')\) 是 \(\mathbf{X}\) 上的滤子基，当且仅当 \(\overline{f}^{1}(\mathbf{M}') \neq \emptyset\)（对每个 \(\mathbf{M}' \in \mathfrak{B}'\)）。这是关系 \(\overline{f}^{1}(\mathbf{M}' \cap \mathbf{N}') = \overline{f}^{1}(\mathbf{M}') \cap \overline{f}^{1}(\mathbf{N}')\) 与第 3 目定义 3 的直接推论。这个条件也可以表述为：\(\mathfrak{B}'\) 的每个集合都与 \(f(\mathbf{X})\) 相交{[}或者说 \(\mathfrak{B}'\) 在 \(f(\mathbf{X})\) 上的迹是一个滤子基{]}。若此条件满足，则 \(f(\overline{f}^{1}(\mathfrak{B}'))\) 是以 \(\mathfrak{B}\) 为基的滤子的一个较细滤子的基。

若 \(\mathfrak{B}\) 是 \(\mathbf{X}\) 上的滤子基，则显然 \(\mathfrak{B}' = f(\mathfrak{B})\) 满足上述条件；于是 \(\overline{f}^1 (f(\mathfrak{B}))\) 是粗于以 \(\mathfrak{B}\) 为基之滤子的某个滤子的基。

设 A 是集合 X 的子集，\(\varphi\) 是典范单射 \(A \rightarrow X\)；若 B 是 X 上的滤子基，则 \(\overline{\varphi}^{1}(\mathfrak{B})\) 与 \(B_{A}\) 相同。若利用上述条件把它表示为 A 上的滤子基，我们就重新得到第 5 目命题 8 的一部分。

